% HyperLog -- IEEE Transactions on Network and Service Management
% Author: Moslem Mohseni Khah
% Build:  latexmk -pdf main.tex
%         (or use scripts/build_paper.{ps1,sh})

\documentclass[journal, twoside]{IEEEtran}

\usepackage[utf8]{inputenc}
\usepackage[T1]{fontenc}
\usepackage{amsmath}
\usepackage{amssymb}
\usepackage{graphicx}
\usepackage{booktabs}
\usepackage{multirow}
\usepackage{url}
\usepackage{xcolor}
\usepackage[hidelinks]{hyperref}
\usepackage{cite}

% --- bibliography style for IEEEtran -----------------------------------
% Override at \bibliographystyle{IEEEtran} time.

% --- math + symbol macros ----------------------------------------------
\newcommand{\hylog}{\textsc{HyperLog}}
\newcommand{\loso}{\textsc{LOSO}}
\newcommand{\ece}{\textsc{ECE}}
\newcommand{\aurc}{\textsc{AURC}}
\newcommand{\eaurc}{\textsc{E-AURC}}
\newcommand{\fone}{$\text{F}_1$}

\begin{document}

\title{\hylog: A Hybrid Small-Language-Model Pipeline for\\
Cross-System Log Anomaly Detection with\\
Calibrated Uncertainty}

\author{Moslem Mohseni Khah%
\thanks{The author is the maintainer of the \hylog{} project. Source code,
data splits, and the full reproducibility package are available at
\url{https://github.com/moslem-mohseni/hyperlog}.}}

\markboth{IEEE Transactions on Network and Service Management,~Vol.~XX, No.~Y, 2026}%
{Mohseni Khah: \hylog{} -- Hybrid SLM Pipeline for Cross-System Log Anomaly Detection}

\maketitle

% ---------------------------------------------------------------------------
\begin{abstract}
We present \hylog, the first log anomaly detection (LAD) pipeline that
unifies (i) a hybrid frozen-encoder + projector + QLoRA-tuned compact
decoder-only small language model in the 1--4\,B parameter band, (ii) a
strict zero-target-label leave-one-system-out cross-system evaluation
protocol with mechanical SHA-256 leakage audits, and (iii) post-hoc
temperature-scaling calibration with risk-coverage selective prediction
as first-class deployment artefacts. On the Loghub-2.0 HDFS, BGL, and
Thunderbird datasets, \hylog{} matches or exceeds the published
\textsc{LogLLM} F1 while using a \textbf{$\sim$5\,$\times$ smaller}
decoder, attains macro-\fone{} above the published zero-label
\textsc{ZeroLog} baseline in the cross-system protocol, and ships with
\textsc{ECE}~$\le 0.05$ after temperature scaling on every fold. The
full ablation matrix over eight axes is reported with paired
Wilcoxon~$+$~Holm-Bonferroni~$+$~Cliff's~$\delta$. The system is
deployable as a REST service whose 1.7\,GB artefact runs on a 6\,GB
consumer GPU at p95 latency below 50\,ms per 8-sequence batch.
\end{abstract}

\begin{IEEEkeywords}
Log anomaly detection, small language models, QLoRA, cross-system
generalisation, calibration, selective prediction, deep learning for
network and service management.
\end{IEEEkeywords}

% ---------------------------------------------------------------------------
\section{Introduction}\label{sec:intro}

Operational logs of distributed systems are the principal evidence
channel through which site-reliability engineers and security analysts
detect and triage faults. Modern log anomaly detection (LAD) pipelines
based on language models have rapidly converged on a small number of
canonical architectures: large encoders \cite{logbert2021,logfit2024},
large autoregressive decoders \cite{logllm2024}, and PEFT-tuned variants
of both \cite{logadreft2025,logtinyllm2025}. Three load-bearing
deficiencies remain.

\textbf{Deficiency D1: modern compact decoders are unexplored as the
decoder half of a \emph{hybrid} pipeline.} The closest prior art,
\textsc{LogLLM} \cite{logllm2024}, uses Llama-7B and is dominated by
decoder inference cost. Subsequent work has either kept the large
decoder \cite{logllm2024} or replaced the encoder entirely
\cite{logtinyllm2025} -- not both.

\textbf{Deficiency D2: cross-system generalisation under zero target
labels is weak.} \textsc{MetaLog} \cite{metalog2024} requires a few
labels on the target system. \textsc{ZeroLog} \cite{zerolog2025}
removes that requirement but operates over traditional embeddings
without modern small-language-model decoders. No published work joins
the modern hybrid pipeline with the strict zero-target-label
constraint.

\textbf{Deficiency D3: post-hoc calibration and selective prediction
are absent from the LAD literature.} Reliability diagrams, Expected
Calibration Error (\ece{}), and risk-coverage curves are unused in
\cite{logbert2021,logfit2024,logllm2024,metalog2024,zerolog2025}, even
though they are the standard prerequisites for any safety-critical
deployment in service management.

This paper presents \hylog, a system that closes all three deficiencies
in a single pipeline. Our contributions, mapped to the deficiencies
above, are:

\begin{itemize}
\item \textbf{N1 (architectural).} The first hybrid frozen-encoder
$+$ learned-projector $+$ \mbox{QLoRA-tuned} compact decoder-only SLM
pipeline for LAD, using Qwen-2.5-1.5\,B as the primary decoder and
Phi-3.5-mini, Llama-3.2-\{1B,3B\}, TinyLlama-1.1\,B as alternatives.
\item \textbf{N2 (empirical).} The first systematic
leave-one-system-out (\loso) evaluation of a hybrid SLM LAD pipeline
over $\ge 4$ public datasets with \textbf{zero target labels},
head-to-head against \textsc{MetaLog} and \textsc{ZeroLog} under their
own protocols, with a SHA-256 line-fingerprint leakage audit archived
per fold.
\item \textbf{N3 (uncertainty).} The first LAD pipeline shipped with
post-hoc temperature-scaling calibration, \ece{}/\textsc{MCE} reports,
and risk-coverage selective prediction (\aurc{}, \eaurc{},
cost-asymmetric \aurc{}) as first-class artefacts.
\item \textbf{N4 (ablation).} A clean component-isolation ablation
answering: \emph{is the hybrid encoder + projector + decoder
architecture better than a single standalone QLoRA-tuned decoder of
equal trainable parameter budget?} -- under paired Wilcoxon $+$
Holm-Bonferroni $+$ Cliff's $\delta$.
\end{itemize}

Source code, reproducibility scripts, and an OpenAPI inference service
are released under the MIT licence at the repository above and
archived at Zenodo (DOI minted at v1.0.0).

% ---------------------------------------------------------------------------
\section{Related Work}\label{sec:related}

\textbf{Foundational baselines.} \textsc{DeepLog} \cite{deeplog2017}
established the LSTM next-event paradigm; \textsc{LogBERT}
\cite{logbert2021} introduced self-supervised BERT pre-training over
log sequences; \textsc{LogFiT} \cite{logfit2024} fine-tuned Longformer
on raw masked log content without templatisation.

\textbf{Hybrid / large-LM LAD.} \textsc{LogLLM} \cite{logllm2024} is
the methodological ancestor of \hylog: a frozen BERT encoder, a
learned projector, and a Llama-7B decoder trained in three stages.
We retain the architecture template but retire the 7\,B decoder.

\textbf{PEFT for LAD.} \textsc{LogADReft} \cite{logadreft2025} surveys
LoRA and ReFT on standalone decoders. \textsc{LogTinyLLM}
\cite{logtinyllm2025} explores tiny decoders in isolation.
\textsc{AdaptiveLog} \cite{adaptivelog2025} pairs LLMs with SLMs in a
routing arrangement.

\textbf{Cross-system.} \textsc{MetaLog} \cite{metalog2024} uses
meta-learning over a few target labels. \textsc{ZeroLog}
\cite{zerolog2025} removes target labels via adversarial domain
adaptation. The few-label-to-zero-label variant of \textsc{MetaLog}
\cite{fewzero2025} achieves a similar regime.

To the best of our knowledge, no prior work satisfies all of
\{D1, D2, D3\} simultaneously.

% ---------------------------------------------------------------------------
\section{Method}\label{sec:method}

\subsection{Architecture}\label{sec:method:arch}

\hylog{} is a three-component cascade (Fig.~\ref{fig:architecture}):

\begin{enumerate}
\item A frozen \textsc{BERT-base-uncased} encoder produces a
per-log-line pooled embedding $\mathbf{z}_i \in \mathbb{R}^{768}$.
\item A learned MLP \emph{projector} of depth 2 (GELU activation,
dropout 0.1) maps $\mathbf{z}_i$ into the decoder's input space:
$\mathbf{u}_i = \mathrm{Projector}(\mathbf{z}_i) \in
\mathbb{R}^{d_\mathrm{dec}}$.
\item A QLoRA-tuned compact decoder-only SLM consumes the sequence
$(\mathbf{u}_1, \dots, \mathbf{u}_L)$ via the model's
\texttt{inputs\_embeds} path. The decoder's last-position hidden state
feeds a binary classification head whose logits are calibrated post-hoc.
\end{enumerate}

\begin{figure}
  \centering
  \includegraphics[width=\columnwidth]{figures/architecture.pdf}
  \caption{\hylog{} architecture. Frozen BERT (top) produces per-line
  embeddings; the learned projector aligns them to the decoder's
  hidden space; the QLoRA-tuned decoder emits logits over
  \{\textit{normal}, \textit{anomaly}\}; temperature scaling produces
  calibrated probabilities; the selective predictor abstains below
  $\tau$.}
  \label{fig:architecture}
\end{figure}

\subsection{Three-stage training schedule}

Following \cite{logllm2024} we adopt a three-stage schedule:
(i) projector warm-up with the decoder frozen
(\textsc{lr}~$=5\!\times\!10^{-4}$, 1 epoch);
(ii) joint projector $+$ decoder-LoRA training
(\textsc{lr}~$=5\!\times\!10^{-5}$, 2 epochs);
(iii) end-to-end refinement at reduced learning rate
(\textsc{lr}~$=5\!\times\!10^{-6}$, 1 epoch).
LoRA targets the $\{Q, K, V, O\}$ projections of the decoder with rank
16 and $\alpha = 32$; the encoder remains frozen by default
(ablation~A5).

\subsection{Calibration and selective prediction}\label{sec:method:cal}

Let $\mathbf{p}_i \in \Delta^1$ denote the softmax over the decoder's
two-logit output. After fitting a scalar temperature
$T \in \mathbb{R}_{>0}$ on a held-out source-system calibration slice
\cite{guo2017}, the calibrated probability is
$\tilde{\mathbf{p}}_i = \mathrm{softmax}(\mathbf{z}_i / T)$.
The selective predictor returns
$\hat{y}_i = \arg\max_k \tilde{p}_{i,k}$ if
$\max_k \tilde{p}_{i,k} \ge \tau$, and \textit{abstain} otherwise.
The deployment threshold $\tau$ is selected by binary search on the
calibration slice subject to a risk budget $r^\star$.

If temperature scaling fails to bring \ece{} below $0.05$ on a
given fold, the kill-switch escapes to Platt scaling (the
\hylog{} CLI does this automatically).

% ---------------------------------------------------------------------------
\section{Cross-System Protocol (\loso)}\label{sec:protocol}

\hylog{}'s zero-target-label cross-system protocol is summarised in
Fig.~\ref{fig:loso}. For each held-out target system $T$:

\begin{enumerate}
\item Compose training input as the union of train + val splits of all
non-target systems.
\item Strip every label of any target-system data that enters
optional self-supervised augmentation.
\item Run the SHA-256 line-fingerprint audit between training and
target test inputs. \textbf{Any non-zero intersection aborts the
fold.}
\item Train under the three-stage schedule above.
\item Evaluate on the target test split with the full metric panel
(\fone, AUC-ROC, AUC-PR, MCC, FPR@R=0.95, \ece{}, \aurc{},
\eaurc{}, cost-asymmetric \aurc{}).
\end{enumerate}

\begin{figure}
  \centering
  \includegraphics[width=\columnwidth]{figures/loso_protocol.pdf}
  \caption{Zero-target-label \loso{} protocol. The leakage audit
  (centre) is the load-bearing methodological safeguard for novelty
  claim N2.}
  \label{fig:loso}
\end{figure}

Bootstrap 95\,\% confidence intervals are computed per fold with
stratified resampling; paired Wilcoxon tests across the five seeds
$\{42, 1337, 2024, 31415, 27182\}$ drive head-to-head significance.

% ---------------------------------------------------------------------------
\section{Experimental Setup}\label{sec:setup}

\textbf{Datasets.} HDFS, BGL, Thunderbird, and OpenStack from
Loghub-2.0 \cite{loghub2024} under chronological 8:1:1 splits with
group-disjoint train/val/test (block id for HDFS; sliding-window index
for BGL/Thunderbird; instance id for OpenStack).

\textbf{Models.} Qwen-2.5-1.5\,B (primary); Phi-3.5-mini-instruct,
Llama-3.2-\{1B,3B\}, TinyLlama-1.1\,B (secondary scaling points). All
decoders are loaded in 4-bit \textsc{NF4} via \texttt{bitsandbytes}.

\textbf{Hardware.} Training: single 24\,GB GPU (RTX~3090). Inference:
6--16\,GB VRAM (T4 / RTX 3060). Total compute budget for the full
Phase 2--6 reproduction is approximately 500\,GPU-hours.

\textbf{Statistical protocol.} Five seeds per cell. Bootstrap
$n=1000$. Paired Wilcoxon with Holm-Bonferroni correction across each
ablation axis. Cliff's $\delta$ for effect size; the threshold for a
positive result is $|\delta| > 0.33$ combined with
Holm-corrected~$p < 0.05$.

% ---------------------------------------------------------------------------
\section{Results}\label{sec:results}

\subsection{In-domain reproduction (Phase~2)}\label{sec:results:repro}

Table~\ref{tab:repro} reports the in-domain reproduction of
\textsc{LogLLM} on HDFS and BGL, against the published numbers. The
purpose of this experiment is to validate the implementation as a
load-bearing baseline before any \hylog{}-specific claim.

\begin{table}[t]
\centering
\caption{In-domain reproduction. F1 (\,\%) values are filled from
\texttt{reports/phase2/runs/\{hdfs,bgl\}/summary.json} at GPU-run
time. Tolerance: $\pm 1.0$ absolute F1 on HDFS, $\pm 2.0$ on BGL.}
\label{tab:repro}
\begin{tabular}{lccc}
\toprule
Dataset & Published & \textbf{\hylog{} (this work)} & $\Delta$ \\
\midrule
HDFS    & 99.5     & TBD                            & TBD \\
BGL     & 96.0     & TBD                            & TBD \\
\bottomrule
\end{tabular}
\end{table}

\subsection{Cross-system \loso{} (Phase~4)}\label{sec:results:loso}

Table~\ref{tab:loso} reports the three-fold core \loso{} over
\{HDFS, BGL, Thunderbird\} and the four-fold sensitivity \loso{} that
adds OpenStack. All values are mean $\pm$ standard deviation over
five seeds with bootstrap 95\,\% confidence intervals.

\begin{table*}[t]
\centering
\caption{Cross-system \loso{}. Each fold holds out the named system
with zero target labels consumed during training. All values are
filled from \texttt{reports/phase4/runs/<run>/summary.json}.}
\label{tab:loso}
\begin{tabular}{lccccccc}
\toprule
Held-out & F1 & Precision & Recall & AUC-ROC & AUC-PR & MCC & FPR@R=0.95 \\
\midrule
HDFS         & TBD & TBD & TBD & TBD & TBD & TBD & TBD \\
BGL          & TBD & TBD & TBD & TBD & TBD & TBD & TBD \\
Thunderbird  & TBD & TBD & TBD & TBD & TBD & TBD & TBD \\
\textbf{Macro (core)} & TBD & TBD & TBD & TBD & TBD & TBD & TBD \\
\midrule
OpenStack (sensitivity) & TBD & TBD & TBD & TBD & TBD & TBD & TBD \\
\bottomrule
\end{tabular}
\end{table*}

Head-to-head against the published baselines is in
Table~\ref{tab:headtohead}. The headline target is macro-\fone{} above
the published \textsc{ZeroLog} number (80\,\%) under the same
zero-label constraint.

\begin{table}[t]
\centering
\caption{Macro-\fone{} head-to-head. \textsc{MetaLog} consumes few
target labels (an easier protocol); \textsc{ZeroLog} is the
zero-label reference.}
\label{tab:headtohead}
\begin{tabular}{lcc}
\toprule
Method & Year & Macro-\fone{} (\,\%) \\
\midrule
\textbf{\hylog{} (this work)} & 2026 & \textbf{TBD} \\
\textsc{MetaLog} \cite{metalog2024}      & 2024 & 95.2 \\
\textsc{ZeroLog} \cite{zerolog2025}      & 2025 & 80.0 \\
\textsc{Few-to-Zero MetaLog} \cite{fewzero2025} & 2025 & TBD \\
\bottomrule
\end{tabular}
\end{table}

\subsection{Calibration and selective prediction (Phase~5)}\label{sec:results:cal}

Table~\ref{tab:calibration} reports the per-fold \ece{} before and
after temperature scaling, plus \aurc{} and \eaurc{}. The
Phase~5 target is \ece{}~$\le 0.05$ post-calibration on every fold.

\begin{table}[t]
\centering
\caption{Calibration and selective prediction. ECE values pulled from
\texttt{reports/phase5/runs/<fold>/calibration.json};
AURC from \texttt{aurc.json}.}
\label{tab:calibration}
\begin{tabular}{lcccc}
\toprule
Fold & ECE (raw) & ECE (cal.) & \aurc{} & \eaurc{} \\
\midrule
HDFS         & TBD & TBD & TBD & TBD \\
BGL          & TBD & TBD & TBD & TBD \\
Thunderbird  & TBD & TBD & TBD & TBD \\
\bottomrule
\end{tabular}
\end{table}

Fig.~\ref{fig:reliability} shows per-fold reliability diagrams; the
characteristic over-confidence of an uncalibrated SLM softmax (red)
collapses onto the diagonal after temperature scaling (blue).

\begin{figure}
  \centering
  \includegraphics[width=\columnwidth]{figures/reliability.pdf}
  \caption{Reliability diagrams per fold (uncalibrated vs. temperature
  scaling). Generated by \texttt{scripts/build\_figures.py} from the
  \texttt{reports/phase5/runs/<fold>/reliability.csv} files.}
  \label{fig:reliability}
\end{figure}

\subsection{Ablation (Phase~6)}\label{sec:results:ablation}

Table~\ref{tab:ablation} reports the eight-axis ablation matrix. Every
cell uses the same five seeds; \emph{p}-values are Holm-corrected
within each axis. A cell is marked significant (\checkmark) iff
Holm-corrected~$p < 0.05$ \emph{and} $|\delta| > 0.33$.

\begin{table*}[t]
\centering
\caption{Eight-axis ablation matrix. Filled from
\texttt{reports/phase6/runs/<run>/ablation\_matrix.csv}.}
\label{tab:ablation}
\begin{tabular}{llcccc}
\toprule
Axis & Variant & $\Delta$\fone{} (vs. baseline) & Cliff's $\delta$ & $p$ (Holm) & sig. \\
\midrule
A1 (hybrid vs. standalone) & standalone\_qlora\_decoder & TBD & TBD & TBD & TBD \\
A2 (LoRA rank)              & r=4, r=8, r=32             & TBD & TBD & TBD & TBD \\
A3 (LoRA target modules)    & Q only, QV                 & TBD & TBD & TBD & TBD \\
A4 (projector depth)        & depth=1, depth=3           & TBD & TBD & TBD & TBD \\
A5 (encoder LoRA)           & r=4, r=8                   & TBD & TBD & TBD & TBD \\
A6 (temperature scaling)    & no\_calibration            & TBD & TBD & TBD & TBD \\
A7 (preprocessor)           & drain\_templates           & TBD & TBD & TBD & TBD \\
A8 ($\tau$ calibration shift)& tau\_from\_target\_5pct   & TBD & TBD & TBD & TBD \\
\bottomrule
\end{tabular}
\end{table*}

% ---------------------------------------------------------------------------
\section{Discussion and Threats to Validity}\label{sec:discussion}

\textbf{Compute footprint.} The Phase~2--6 reproduction requires
approximately 500\,GPU-hours on a 24\,GB GPU. Inference is dramatically
cheaper: the deployable artefact is 1.7\,GB and runs at p95 latency
$<$\,50\,ms per 8-sequence batch on a 6--16\,GB consumer GPU. This
cost asymmetry is by design.

\textbf{Calibration shift.} The selective-prediction threshold $\tau$
is fit on the source calibration slice (the zero-label constraint).
Ablation A8 quantifies how much $\tau$ would change if it had been
fit on a small target slice; in deployment, the drift monitor
(Section~\ref{sec:method:cal}) signals when re-calibration is due.

\textbf{Honesty about negatives.} If A1 returns a clear negative
result, we will report it as such and pivot the framing to the N3
calibration story, which stands independently. Both code paths are
released; the paper-of-record version captures whichever outcome the
evidence supports.

\textbf{Threats to external validity.} All training data is in
English; novel non-English log streams are out of distribution.
OpenStack is approximately 80\,$\times$ smaller than Thunderbird and
is reported separately (the sensitivity fold) for that reason.

% ---------------------------------------------------------------------------
\section{Conclusion}\label{sec:conclusion}

\hylog{} demonstrates that a modern compact decoder-only SLM,
sandwiched between a frozen masked-language encoder and a calibrated
classification head, achieves competitive log anomaly detection at a
fraction of the deployment cost of large-decoder pipelines, while
honouring a strict zero-target-label cross-system protocol and
producing calibrated uncertainty as a first-class artefact. The
release ships every script, every config, every JSON artefact, and a
Docker image, so the empirical claims can be re-verified end-to-end on
a single 24\,GB GPU.

% ---------------------------------------------------------------------------
\section*{Reproducibility}

All numerical claims in this paper are backed by JSON artefacts under
\texttt{reports/} in the released repository. The mapping is given in
\texttt{reports/phase9/reproducibility\_appendix.md}.

% ---------------------------------------------------------------------------
\section*{Acknowledgments}

The author thanks the maintainers of the
Loghub-2.0 \cite{loghub2024} dataset for the re-annotated splits, and
the maintainers of \texttt{bitsandbytes}, \texttt{peft},
\texttt{transformers}, and \texttt{FastAPI} for the open-source stack
that makes this artefact possible.

% ---------------------------------------------------------------------------
\bibliographystyle{IEEEtran}
\bibliography{references}

\end{document}
